\section{Conclusion}
\label{sec:conclusion}

This work introduced a reference-aware evaluation framework for neural wavefield emulators, isolating same-reference surrogate error, raw solver discrepancy, and cross-reference surrogate discrepancy across Hydrostatic, MUSCL-HR, and Boussinesq references. Aligning outputs to 50 shared requested output times on the central $64 \times 64$ crop of a buffered $192 \times 192$ computation prevents temporal sampling mismatch from being conflated with model error without assuming an absolute physical ground truth.

The evaluation shows global solver-gap RMSEs of $0.00171$, $0.00305$, and $0.00374$ for the Hydrostatic--MUSCL-HR, MUSCL-HR--Boussinesq, and Hydrostatic--Boussinesq pairs. These gaps are smaller than some direct surrogate errors and comparable to or larger than others, depending on the target and metric. ConvLSTM is the strongest selected seed-$18$ Hydrostatic direct model by global relative $L_2$, followed by U-FNO and F-FNO, while all six cross-reference ratios exceed one. The ratios are benchmark-relative effects rather than evidence of physical generalization beyond the training references. Wet--dry stress failures, family-specific holdout degradation, and under-dispersed uncertainty bounds further mark the scope; the study makes no operational latency claim.

Future developments require physical-scale surveyed bathymetry, realistic fault rupture dynamics, wet--dry coastal boundary conditions, and direct assimilation of gauge observations. By establishing a controlled, transparent benchmarking framework, this study clarifies how conclusions regarding neural PDE surrogates depend directly on underlying numerical references and temporal alignment procedures.
